% 2.1 Grundlegende Begriffsdefinitionen
% - Definition der Schlüsselbegriffe aus dem Kontext der Arbeit, die für das
% weitere Verständnis notwendig sind
% - Allgemein technische Vorkenntnisse des Lesers voraussetzen!
% - Begriffsdefinitionen werden wörtlich zitiert („ ...“) mit exakter Quellenangabe
% (siehe Abschnitt Quellenarbeit)
% - ggf. Ableitung eigener Begriffsdefinitionen
\subsection{Begriffliche Grundlagen und Datenqualität}

\subsection{Technologische Grundlagen}

% 2.2 – 2.n-1 Kurze Erläuterung von n-1 bestehenden Ansätzen
% - Konkurrierende Ansätze: Ansätze, die in Konkurrenz zum eigenen Thema
% stehen
% - Komplementäre Ansätze: Ansätze, die in der Arbeit verwendet werden / auf
% denen aufgesetzt wird
% - weitere Ansätze aus dem Themenfeld (z. B. aus der Praxis)
\subsection{Verwandte Arbeiten und Identifikation der Forschungslücken}